\documentclass[10pt,a4paper]{article}
\usepackage{rivendell}
% Options: \usepackage[times]{rivendell}    Times fonts
%          \usepackage[numbers]{rivendell}  numeric citations
%          \usepackage[nonatbib]{rivendell} bring your own biblatex
%          \usepackage[draft]{rivendell}    DRAFT watermark + visible \todo notes
%          \usepackage[trackchanges]{rivendell}  show \added/\deleted/\replaced

% \course{[CS7.401] Introduction to Something}
\assignment{Assignment 1}
\shorttitle{Example Report}   % optional, right side of the running head
% \headertext{Just a Report or Whatever}  % or set the head freely

\title{An Example Report Typeset with \texttt{rivendell}\thanks{Submitted on \today.}}

\author[1]{First Author}
\author[1,2]{Second Author}
\affil[1]{International Institute of Information Technology, Hyderabad}
\affil[2]{Some Other Lab, Somewhere}
\email{first.author@example.com}
\email{second.author@example.com}
% \collaboration{(Team 7)}
% \date{}   % hide the date

\begin{document}

\maketitle

\begin{abstract}
This document shows every feature of the \texttt{rivendell} package, which gives the
standard \texttt{article} class the look of the AAS journals' default style. Drop
\texttt{rivendell.sty} (and \texttt{rivendell.bst} if you cite from a \texttt{.bib} file)
next to your \texttt{.tex} file and load it with \verb|\usepackage{rivendell}|.
\end{abstract}

\keywords{Templates; Typesetting; Assignments}

\section{Introduction}\label{sec:intro}

Section headings are centred and set in capitals; subsection headings are centred
italics. Citations are author--year by default, for example \citet{lamport94} or
\citep{knuth84}. Cross-references made with \verb|\cref| name themselves:
\cref{sec:method}, \cref{fig:example}, \cref{tab:results}, \cref{eq:ridge},
\cref{lst:ridge} and \cref{app:extra}.\footnote{Ordinary
footnotes work as usual.}

\subsection{Background}

Let $\mathbf{X} \in \mathbb{R}^{n \times d}$ be a design matrix and
$\mathbf{y} \in \mathbb{R}^{n}$ a target vector. For $\lambda > 0$ the ridge-regression
estimate is
\begin{equation}
  \hat{\mathbf{w}} = (\mathbf{X}^\top \mathbf{X} + \lambda \mathbf{I})^{-1}
  \mathbf{X}^\top \mathbf{y},
  \label{eq:ridge}
\end{equation}
where $\mathbf{I}$ is the $d \times d$ identity matrix.

\subsubsection{A third-level heading}

Third-level headings are smaller still.

\paragraph{Run-in heading} Paragraph headings run into the text.

\section{Method}\label{sec:method}

\clearpage

\begin{figure}[t]
  \centering
  \includegraphics[width=0.45\linewidth]{example-image}
  \caption{Figure captions get a bold label and sit below the figure.}
  \label{fig:example}
\end{figure}

\begin{table}[t]
  \centering
  \caption{Table captions sit above the table. Rules come from \texttt{booktabs}.}
  \label{tab:results}
  \begin{tabular}{lcc}
    \toprule
    Method & Accuracy & Time (s) \\
    \midrule
    Baseline & 0.81 & 12.4 \\
    Ours     & 0.87 & 10.9 \\
    \bottomrule
  \end{tabular}
  \tablenote[0.5\linewidth]{Numbers are made up for illustration.}
\end{table}

Some text about the method, referring back to \cref{eq:ridge}. \Cref{lst:ridge}
implements it with \lstinline|np.linalg.solve| rather than an explicit inverse.

\begin{lstlisting}[language=Python, caption={Ridge regression in NumPy.}, label=lst:ridge]
import numpy as np

def ridge(X, y, lam):
    """Closed-form ridge estimate."""
    d = X.shape[1]
    return np.linalg.solve(X.T @ X + lam * np.eye(d), X.T @ y)  # (d,)
\end{lstlisting}

With the \texttt{trackchanges} option, edits are marked: this sentence was
\replaced{rewritten}{written badly}, \added{some words were added,} and some
\deleted{unnecessary }words were removed. Without the option only the final text is
printed. \todo{Notes like this one appear only with the \texttt{draft} option.}

\begin{acknowledgments}
We thank the teaching assistants.
\end{acknowledgments}

\appendix

\section{Extra derivations}\label{app:extra}

Appendix sections are lettered, and their equations are numbered accordingly:
\begin{equation}
  \lVert \mathbf{x} \rVert_2^2 = \sum_{i=1}^{d} x_i^2 .
\end{equation}

\bibliographystyle{rivendell}
\begin{thebibliography}{2}
\expandafter\ifx\csname natexlab\endcsname\relax\def\natexlab#1{#1}\fi
\bibitem[{Knuth(1984)}]{knuth84}
Knuth, D. E. 1984, The \TeX book (Addison-Wesley)
\bibitem[{Lamport(1994)}]{lamport94}
Lamport, L. 1994, \LaTeX: A Document Preparation System, 2nd edn.\ (Addison-Wesley)
\end{thebibliography}
% With a .bib file, replace the thebibliography block above with:
% \bibliography{references}

\end{document}
